\section{Ratios de liquidez}
Los ratios de liquidez evalúan la capacidad de atender obligaciones de corto plazo. Su interpretación exige revisar la calidad de los activos corrientes: no todo activo corriente se convierte en efectivo con la misma rapidez.

\subsection{Razón corriente}
\textbf{Fórmula:}
\[
\text{Razón corriente}=\frac{\text{Activo corriente}}{\text{Pasivo corriente}}
\]
\textbf{Caso académico elaborado con fines educativos.} Comercial Andina S.A.C. presenta activo corriente de \soles{480,000.00} y pasivo corriente de \soles{300,000.00}.
\[
\frac{480,000.00}{300,000.00}=1.60
\]
\textbf{Interpretación:} por cada \soles{1.00} de obligación corriente, la empresa dispone de \soles{1.60} en activos corrientes. No debe afirmarse automáticamente que el resultado es favorable sin revisar inventarios, cobrabilidad y sector.

\textbf{Decisión:} revisar la estructura del activo corriente y programar pagos según la disponibilidad real de efectivo.

\newpage
\subsection{Prueba ácida}
\textbf{Fórmula:}
\[
\text{Prueba ácida}=\frac{\text{Activo corriente}-\text{Inventarios}}{\text{Pasivo corriente}}
\]
Datos: activo corriente \soles{480,000.00}, inventarios \soles{180,000.00} y pasivo corriente \soles{300,000.00}.
\[
\frac{480,000.00-180,000.00}{300,000.00}=1.00
\]
\textbf{Interpretación:} sin depender de la venta de inventarios, la empresa posee \soles{1.00} de activos líquidos por cada \soles{1.00} de deuda corriente.

\textbf{Decisión:} fortalecer caja y cobranza, porque el margen ajustado es exacto.
